% -----------------------------------------------------------------------
% REWRITTEN conclusion — expanded to ~12 lines; NLP thesis bookend;
% future work; dataset release; \expnote for pending items
% -----------------------------------------------------------------------

\section{Conclusion}

Question linguistic structure is a principled and practical signal for pre-retrieval fidelity selection: \emph{how many}, \emph{what color}, and \emph{read the text} reliably require high-resolution evidence, while \emph{is there} and \emph{what type} are resolvable from low-fidelity representations.
\textsc{VOILA} operationalizes this insight via a gradient-boosted classifier, isotonic calibration, and a greedy value-of-information decision rule---selecting the minimum-cost sufficient fidelity from question text alone, before any image is accessed---reducing retrieval cost by 50--65\% across five datasets and six VLMs (7B--235B parameters) while matching or exceeding full-resolution accuracy on OCR-heavy and domain-specialized benchmarks, and remaining within an average of 10 accuracy points on general VQA.
Natural extensions include richer question encoders for ambiguous queries, multilingual support, and fidelity selection for video (frame rate) and document retrieval (chunk size); code, checkpoints, and the LoCoMo-VQA benchmark will be released.
\fablenote{For AAAI, close on the decision-theoretic thesis rather than the linguistic one: e.g., ``More broadly, VOILA demonstrates that the value of information an AI system acquires can be estimated before paying for it --- a principle that extends beyond images to any perception pipeline with tiered acquisition costs.'' This bookends the AAAI framing suggested for the abstract and positions the paper as a general contribution to resource-bounded AI rather than a VQA-specific trick. Also reconcile ``50--65\%'' here with whichever range the abstract settles on.}
